\documentclass[border=3mm]{standalone}
\usepackage{tikz}
\usepackage{xeCJK}
\setCJKmainfont{Noto Sans CJK SC}
\usetikzlibrary{positioning,arrows.meta,fit,backgrounds,calc,
                 decorations.pathreplacing}
\usepackage{amsmath,amssymb}
\begin{document}
\begin{tikzpicture}[font=\small,node distance=4mm and 5mm]

% ========== 配色 ==========
\definecolor{cB}{HTML}{4A90D9}   % 蓝-节点
\definecolor{cO}{HTML}{E8913A}   % 橙-边/距离
\definecolor{cG}{HTML}{5BAA6A}   % 绿-backbone
\definecolor{cH}{HTML}{E8913A}   % 橙-任务头
\definecolor{cK}{HTML}{555555}   % 深灰

% ========== 样式 ==========
\tikzset{
  arr/.style={-{Latex[length=1.6mm,width=1.2mm]},line width=.55pt,cK},
  blk/.style={rounded corners=2.5pt,draw=#1,line width=.6pt,
    fill=#1!8,align=center,inner sep=3pt,font=\footnotesize},
  blk/.default=cK,
  bblk/.style={rounded corners=3pt,draw=cG,line width=.7pt,
    fill=cG!6,align=center,inner sep=3pt,font=\footnotesize},
  hblk/.style={rounded corners=2.5pt,draw=cH,line width=.6pt,
    fill=cH!8,align=center,inner sep=3pt,font=\footnotesize},
  panel/.style={rounded corners=4pt,draw=cK!25,line width=.4pt,
    fill=black!2,inner sep=5pt},
  plabel/.style={font=\footnotesize\bfseries,anchor=north west,text=cK},
}

% ================================================================
%  (a) 图表示
% ================================================================
\begin{scope}[shift={(0,0)}]
  % 配体（六边形 + 原子）
  \coordinate(L0)at(0,0);
  \draw[thick,cB!70,fill=cB!8,rounded corners=1pt]
    (-.55,.32)--(.1,.65)--(.75,.32)--(.75,-.32)--(.1,-.65)--(-.55,-.32)--cycle;
  \foreach \x/\y in {-.55/.32,.1/.65,.75/.32,.75/-.32,.1/-.65,-.55/-.32}{
    \fill[cB] (\x,\y) circle(1.5pt);}
  \node[font=\scriptsize,below=1.5mm,text=cB]at(0.1,-.65){配体};

  % 口袋（弧形 + 残基）
  \begin{scope}[shift={(2.4,0)}]
    \draw[thick,cO!55,fill=cO!5,rounded corners=7pt]
      (-.45,.8)--(0,1)--(.65,.8)--(.8,.2)--(.65,-.45)--(0,-.7)--(-.45,-.45)--(-.6,.2)--cycle;
    \foreach \x/\y in {-.25/.55,.3/.65,.55,.25,.45/-.15,.1/-.4,-.35/-.05}{
      \fill[cO] (\x,\y) circle(1.5pt);}
    \node[font=\scriptsize,below=1.5mm,text=cO]at(0,-.7){口袋};
  \end{scope}

  \draw[arr](3.5,0)--(4.3,0);

  % ---- 节点特征（竖向条形 = 每个原子一个向量）----
  \begin{scope}[shift={(4.6,.5)}]
    \foreach \i in {0,...,5}{
      \pgfmathsetmacro{\h}{0.5+rnd*0.5}
      \fill[cB!50] (\i*2.8mm,0) rectangle ++(2.4mm,-\h cm);
    }
    \draw[cK!35,line width=.3pt](0,.05) rectangle (16.8mm,-1.1cm);
    \node[font=\tiny,anchor=north,text=cK]at(8.4mm,-1.15cm){节点特征};
    \node[font=\tiny,anchor=south,text=cB!70]at(8.4mm,.1){每原子一维向量};
  \end{scope}

  % ---- 边特征（对称热力图 = 原子对距离矩阵）----
  \begin{scope}[shift={(7.0,.5)}]
    % 画成对角对称的小方块（距离矩阵感觉）
    \foreach \r in {0,...,4}{
      \foreach \c in {0,...,4}{
        \pgfmathsetmacro{\dist}{abs(\r-\c)}
        \pgfmathsetmacro{\val}{int(max(10, 80-\dist*18+rnd*12))}
        \fill[cO!\val] (\c*2.3mm,-\r*2.3mm) rectangle ++(2mm,2mm);
      }
    }
    % 对角线高亮
    \foreach \i in {0,...,4}{
      \fill[cO!90] (\i*2.3mm,-\i*2.3mm) rectangle ++(2mm,2mm);
    }
    \draw[cK!35,line width=.3pt](0,-11.5mm) rectangle (11.5mm,2mm);
    \node[font=\tiny,anchor=north,text=cK]at(5.75mm,-12mm){边特征};
    \node[font=\tiny,anchor=south,text=cO!70]at(5.75mm,2.5mm){原子对距离矩阵};
  \end{scope}
\end{scope}

% (a) 面板
\coordinate(pa_tl)at(-1.3,1.5);
\coordinate(pa_br)at(9.5,-1.5);
\node[plabel]at(pa_tl){\textsf{(a) 图表示}};
\begin{scope}[on background layer]
  \node[panel,fit=(pa_tl)(pa_br)](pa){};
\end{scope}

% ================================================================
%  (b) 共享编码主干
% ================================================================
\begin{scope}[shift={(0,-3.6)}]

  \node[plabel]at(-1.1,1.1){\textsf{(b) 共享编码主干}};

  % 编码模块
  \node[bblk,minimum width=18mm,minimum height=11mm]at(1.0,0)(intra){%
    分子内\\编码器$\times\!12$};
  \node[bblk,minimum width=18mm,minimum height=11mm,right=6mm of intra](inter){%
    分子间\\编码器$\times\!6$};
  \node[bblk,minimum width=18mm,minimum height=11mm,right=6mm of inter](eol){%
    边输出层};

  \draw[arr]([xshift=-7mm]intra.west)--(intra.west);
  \draw[arr](intra)--(inter);
  \draw[arr](inter)--(eol);
  \draw[arr](eol.east)--++(7mm,0);

  % 标注：参数共享
  \node[font=\tiny,text=cG!70,anchor=south]
    at([yshift=2mm]inter.north){所有任务共享同一组参数};

  % 面板
  \coordinate(pb_tl)at(-1.3,1.3);
  \coordinate(pb_br)at(8.3,-1.0);
  \begin{scope}[on background layer]
    \node[panel,fit=(pb_tl)(pb_br)](pb){};
  \end{scope}
\end{scope}

% ================================================================
%  (c) 多任务输出头
% ================================================================
\begin{scope}[shift={(0,-6.6)}]

  \node[plabel]at(-1.1,1.1){\textsf{(c) 多任务输出头}};

  % 虚拟节点
  \node[blk=cG,minimum width=13mm,minimum height=8mm]at(0.5,.0)(vn){%
    虚拟\\[-1pt]节点};

  % 三个头
  \node[hblk,minimum width=16mm,minimum height=8mm,
    right=8mm of vn,yshift=9mm](haff){%
    {\scriptsize\textbf{亲和力回归}}\\[-1pt]{\tiny FFN$\to\mathrm{p}K_d$}};
  \node[hblk,minimum width=16mm,minimum height=8mm,
    right=8mm of vn](hpose){%
    {\scriptsize\textbf{姿态选择}}\\[-1pt]{\tiny FFN$\to$logit}};
  \node[hblk,minimum width=16mm,minimum height=8mm,
    right=8mm of vn,yshift=-9mm](henergy){%
    {\scriptsize\textbf{能量预测}}\\[-1pt]{\tiny MDN+原子}};

  \draw[arr](vn.east)--(haff.west);
  \draw[arr](vn.east)--(hpose.west);
  \draw[arr](vn.east)--(henergy.west);

  % task_mask 标注
  \node[font=\scriptsize,text=cH,right=3mm of haff]{亲和力批次激活};
  \node[font=\scriptsize,text=cH,right=3mm of hpose]{亲和力批次激活};
  \node[font=\scriptsize,text=cB,right=3mm of henergy]{能量批次激活};

  % 面板
  \coordinate(pc_tl)at(-1.3,1.3);
  \coordinate(pc_br)at(8.3,-1.5);
  \begin{scope}[on background layer]
    \node[panel,fit=(pc_tl)(pc_br)](pc){};
  \end{scope}
\end{scope}

% ================================================================
%  面板间箭头
% ================================================================
\draw[arr,thick]([yshift=-2mm]pa.south)--([yshift=2mm]pb.north)
  node[midway,left,font=\tiny]{特征输入};
\draw[arr,thick]([yshift=-2mm]pb.south)--([yshift=2mm]pc.north)
  node[midway,left,font=\tiny]{共享表示};

\end{tikzpicture}
\end{document}
